\section{Limits and Deployment Implications}\label{sec:discussion}
\paragraph{What the results establish.}
The replay separates three kinds of evidence. Zero policy and headroom
violations show that selected nodes obey the corresponding checks.
Model metrics measure prediction on held-out recorded transitions.
Migration counts and the next-state proxy describe how those checks change
decisions. The first is largely an invariant of filtering; it should not
be read as independent evidence of a model's accuracy or production safety.

The strongest reported comparative result is lower migration frequency
than the current-state learned control. The natural-workload next-state
difference against that control has a confidence interval reaching zero.
Adding Q95 to the temporal scheduler removes violations of the Q95 rule
but does not improve the reported next-state proxy. Its measured benefit
is compliance with that additional rule, at the cost of more migration
decisions. Neither finding supports an unconditional claim of superiority.

\paragraph{Missing feedback and resource accounting.}
A selected destination does not receive additional load in this replay,
and the source is not relieved of it. Repeated decisions therefore do not
accumulate the physical effects that a live controller would produce.
Instance signals trigger events but are not translated into destination
CPU or memory demand. Before deployment, a controller needs resource
requests, capacity reservations, migration costs, and a model of transient
resource use. The current headroom rule cannot replace that accounting.

\paragraph{Scope of generalization.}
The evaluation uses one hour-long microservice shard from one provider,
with a test interval of 59 timestamps. Synthetic policy overlays probe
variation within one generator; they do not reproduce real geographic,
trust, or service-sensitivity distributions. The balanced workload shares
the same telemetry as the natural workload. The reported bootstrap also
does not retain serial dependence between timestamps. Additional temporal
windows and a block-bootstrap sensitivity analysis would strengthen the
uncertainty assessment.

Destination fairness depends on both the number of migrations and the
eligible sets. Jain values across workloads with different event counts
are not directly comparable. History $H$ counts retained sources as well
as migrations, whereas the reported fairness metric counts migrations
only. Neither captures tenant entitlement or heterogeneous node capacity.

\paragraph{A possible controller integration.}
A Kubernetes implementation could use a Filter plugin for policy and
admission predicates, shared state for telemetry and history, and
Reserve/Unreserve plus a final pre-bind check for coordinated decisions
\citep{kubernetes2026}. Reproducing the shortlist's lexicographic selection
requires explicit coordination with scoring; an ordinary weighted Score
plugin does not by itself implement this ordering. Moving existing
instances also requires a separate reconciliation and migration mechanism.
This is a proposed mapping, not an implemented or measured integration.

Policy revocation must invalidate cached masks. Stale telemetry, model
drift, controller restarts, and concurrent schedulers require explicit
handling. A failed recheck must leave the event uncommitted and release
any tentative reservation. Testing these paths, measuring full control-plane
latency, and evaluating post-placement application behavior remain necessary
before making deployment claims.
